\section*{Appendix}
\section{Sequential Meta-Analysis}
\label{sec:Meta_analysis}
Here we want to briefly give a summary of the fixed-effect and random-effects models used in (sequential) meta-analysis. 
Classical \citep{Whitehead_2002,Cooper_2009} and cumulative \citep{Higgins_Whitehead_Simmonds_2011} meta-analysis are concerned with the 
aggregation of (unbiased) effect estimates. In this setup, every study reports an effect estimate $\hat \vartheta_i$ and corresponding standard error estimate $\hat \sigma_i$.

Fixed-effect meta-analysis\footnote{One can distinguish between the common effect and fixed-effects (plural) analysis in which the study effects are unknown and fixed but need not necessarily be identical. While the interpretation of $\vartheta$ changes, the estimators remain identical \citep{Rice_2018}.} builds on the assumption that the effect is constant across studies. In this case, the following estimates and standard error for the aggregate effect estimate 
\begin{align}
\label{eq:FE}
    \hat \vartheta_{\text{FE}} = \frac{\sum_{i=1}^n \hat \vartheta_i/\hat \sigma_i^2}{\sum_{i=1}^n 1/\hat \sigma_i^2}
\hspace{1cm} \text{and}\hspace{1cm} 
   \widehat{\text{Var}} (\hat \vartheta_{\text{FE}}) = \frac{1}{\sum_{i=1}^n 1/\hat \sigma_i^2}
\end{align}
are widely used. This estimator can be understood as weighted-least squares (with the inverse variance as weights) taking the variance estimates of the individual studies at face value (i.e., assuming that the variance of the variance estimate is negligible).
Alternatively, and in line with the Bayesian perspective taken in this paper, we can understand the fixed-effect model as 
\begin{align}\nonumber
    \hat \vartheta_i \lvert \vartheta \sim \text{normal}(\vartheta, \sigma^2_i).
\end{align}
The posterior of the fixed effect (assuming conditional independence) is given by
\begin{align}\nonumber
    p(\vartheta \lvert \hat \vartheta_1, \dots , \hat \vartheta_n) = \frac{\left(\prod_{i=1}^n \text{normal}(\vartheta, \sigma^2_i)\right) \pi(\vartheta)}{p(\hat \vartheta_1, \dots , \hat \vartheta_n)}.
\end{align}
If the prior $\pi(\vartheta)$ is chosen to be flat, the posterior mean $\mathbb{E}[\vartheta \lvert \hat \vartheta_1, \dots , \hat \vartheta_n] = \hat \vartheta_{FE}$ is equivalent to the classical fixed-effects estimator in Equation \ref{eq:FE} (again assuming the variance of the variance estimates to be negligible, i.e. assuming $\sigma_i$ to be known).

While effect homogeneity may be a reasonable assumption in some settings, it is doubtful in biomedical settings or the social sciences. Studies vary in numerous ways including their target population, their intervention, and their outcome measures. For these reasons, at least some degree of effect heterogeneity is to be expected \citep{Higgins_Thompson_Spiegelhalter_2009}. 
Following this line of thought, random-effects models recommend themselves.\footnote{There are various interpretations of random effects models \citep{Higgins_Thompson_Spiegelhalter_2009,Cooper_2009,Gelman_2005}. While we prefer the Bayesian interpretation below, random effects are also frequently understood by reference to a sampling mechanism by which effects for every study are sampled from a population of possible effects.} 
If we assume that the effects $\vartheta_i$ of the different studies are exchangeable - that is, they may be non-identical but we have no reason to distinguish their magnitude a priori - then de Finetti's theorem allows us to think of the $\vartheta_i$ as identically and independently distributed conditional on some (prior) parameter $\psi$ \citep{Bernardo_1994}.\footnote{More specifically, the sequence ${\vartheta_i}$ needs not only to be exchangeable but embeddable in an infinite sequence that is exchangeable.} 
This motivates the random-effects model \citep{Whitehead_2002}.

\begin{align}\nonumber
    \hat \vartheta_i \lvert \mu_i\, & \sim \,\text{normal}(\mu_i , \sigma^2_i) \\\nonumber
    \mu_i \lvert \vartheta,\tau \, & \sim \,\text{normal}(\vartheta,\tau^2), 
\end{align}
where $\vartheta$ is the mean of the study effects and $\tau$ describes the level of effect heterogeneity.
Integrating out $\mu_i$, we have 
\begin{align}\nonumber
    \hat \vartheta_i \lvert \vartheta, \tau \, & \sim \,\text{normal}(\vartheta, \sigma_i^2+\tau^2). 
\end{align}
The posterior of the mean of the study effects is then given by (assuming $\vartheta$ is independent of $\tau$)
\begin{align}\nonumber
     p(\vartheta \lvert \hat \vartheta_1, \dots , \hat \vartheta_n,\tau) = \frac{\left(\prod_{i=1}^n \text{normal}(\vartheta, \sigma_i^2+\tau^2)\right) \pi(\vartheta)}{p(\hat \vartheta_1, \dots , \hat \vartheta_n,\tau)}.
\end{align}
Choosing a flat prior for $\vartheta$ (assuming an appropriate plug-in estimator $\hat \tau$ for $\tau$ is available and neglecting the variance of $\hat \sigma_i$) yields the typically used random-effects estimator \citep{Whitehead_Whitehead_1991,Higgins_Thompson_Spiegelhalter_2009}:
\begin{align}
\label{eq:REest}
    \hat \vartheta_{RE} = \frac{\sum_{i=1}^n\frac{1}{\hat\sigma_i^2 + \hat \tau^2 }\hat \vartheta_i}{\sum_{i=1}^n\frac{1}{\hat\sigma_i^2 + \hat \tau^2 }} \hspace{0.3cm} \text{and}\hspace{0.3cm} 
   \widehat{\text{Var}} (\hat \vartheta_{RE}) = \frac{1}{\sum_{i=1}^n\frac{1}{\hat\sigma_i^2 + \hat \tau^2 }}.
\end{align}

From the discussion above, it is clear what the fixed- and random-effects estimates are when these models are applied sequentially. In this case we understand $\{(\hat \vartheta_i,\hat \sigma_i)\}_{i=1}^n$ as a sequence that we observe up to point  $m \leq n$. To perform meta-analysis sequentially, we apply the above models to $\{(\hat \vartheta_i,\hat \sigma_i)\}_{i=1}^m$ for every $1\leq m \leq n$.\footnote{As discussed in Section \ref{Sec:QuantifyingLearning}\ref{sec:random-effects-model}, the posterior of the random-effects model can also be updated in fully Bayesian fashion (that is, without the use of a plug-in estimator) using modern probabilistic programming languages, such as \citep{stan}}.

This  reveals that the fixed-effects model is
not conceptually suited to explain the kind of increased uncertainty observed in Section \ref{Sec:Studies}\ref{Sec:ToyExample}. As can be seen from Equation \ref{eq:FE}, when performing fixed-effects meta-analysis, every future study will increase our certainty (regardless of how surprising its outcome may be). 

For random-effect models, uncertainty may increase if a future study increases our estimate of effect heterogeneity $\hat \tau$  (see Equation \ref{eq:REest}). Classically, this increase in uncertainty, however, is solely attributed to effect heterogeneity. In Section \ref{Sec:QuantifyingLearning}\ref{sec:random-effects-model}, we offer an alternative interpretation of the random-effects model in terms of methodological biases. This allows us to understand the increase in posterior variance when the posterior is updated sequentially as a consequence of a new study forcing us to question our beliefs about the corpus of studies collected so far.  

The example in Section \ref{Sec:Studies}\ref{Sec:ToyExample} illustrates one case in which uncertainty may increase through a process other than effect heterogeneity.  In this example, revelations about systematic measurement error lead to a sudden change in posteriors.
While measurement error has been discussed in the meta-analysis literature, 
it is usually assumed to only attenuate effect estimates \citep{Cooper_2009}. Recent work, however, highlights that correlated measurement error can systematically bias effect estimates \citep{Knox_Lucas_Cho_2022,Mikhaeil_2025}. The wider issue of bias has been addressed in the meta-analysis literature \citep{Higgins_Thompson_Spiegelhalter_2009,Carter_2019,Wiernik_Dahlke_2020}, and both careful assessment of bias and selection of studies to be included are common. 


\section{Comparison of Meta-Analysis Weights with our Sequential Meta-Analysis Trace}
\label{Sec:CompMAW_SMAT}
\begin{table}\renewcommand\thetable{S1} 
\centering
\begin{tabular}[t]{c|cc|cc}
Number in  & Sequential & Sequential  & Retrospective & Retrospective\\
Sequence & FE weight & RE weight & FE weight & RE weight\\
\midrule
1 & 100.0 & 100.0 & 3.3 & 3.3\\
2 & 49.4 & 49.4 & 3.2 & 3.3\\
3 & 39.7 & 39.7 & 4.3 & 3.4\\
4 & 19.8 & 19.8 & 2.7 & 3.3\\
5 & 22.9 & 22.9 & 4.0 & 3.4\\
6 & 17.1 & 17.1 & 3.6 & 3.3\\
7 & 11.6 & 11.6 & 2.8 & 3.3\\
8 & 8.8 & 8.8 & 2.3 & 3.3\\
9 & 11.6 & 11.6 & 3.4 & 3.3\\
10 & 11.8 & 11.8 & 4.0 & 3.4\\
11 & 7.9 & 9.0 & 2.9 & 3.3\\
12 & 7.1 & 8.2 & 2.8 & 3.3\\
13 & 12.6 & 7.9 & 5.7 & 3.4\\
14 & 6.2 & 7.1 & 3.0 & 3.3\\
15 & 6.9 & 6.7 & 3.5 & 3.3\\
16 & 5.6 & 6.2 & 3.1 & 3.3\\
17 & 6.1 & 5.9 & 3.5 & 3.3\\
18 & 6.8 & 5.6 & 4.2 & 3.4\\
19 & 4.1 & 5.2 & 2.6 & 3.3\\
20 & 4.2 & 5.0 & 2.9 & 3.3\\
21 & 4.7 & 4.8 & 3.3 & 3.3\\
22 & 4.6 & 4.6 & 3.4 & 3.3\\
23 & 4.3 & 4.4 & 3.3 & 3.3\\
24 & 4.6 & 4.2 & 3.7 & 3.4\\
25 & 3.3 & 4.0 & 2.8 & 3.3\\
26 & 3.7 & 3.8 & 3.2 & 3.3\\
27 & 2.9 & 3.7 & 2.6 & 3.3\\
28 & 3.7 & 3.6 & 3.5 & 3.3\\
29 & 3.3 & 3.4 & 3.2 & 3.3\\
30 & 3.2 & 3.3 & 3.2 & 3.3\\
\end{tabular}
\caption{\label{tab:MAweights} Fixed-effect (FE) and random-effects (RE) meta-analysis weights for the stylized example in Section \ref{Sec:Studies}\ref{Sec:ToyExample}. For the sequential meta-analysis weights, meta-analysis is performed sequentially and the weight of the most recent study is reported. The retrospective weights are obtained by performing meta-analysis on the entire sequence of studies. Meta-analysis weights, either computed sequentially or retrospectively, are not able to capture the influential nature of 11th study. }
\end{table}

Our sequential meta-analysis method to quantify how much has been learned is complementary to the classical meta-analysis perspective. Through the lens of meta-analysis, the value of a study reflects its contribution to the overall meta-analysis effect estimate. This view judges a study in comparison to all other studies as a function of each study's estimated standard error and the overall degree of effect heterogeneity. No account is taken of whether certain studies may have been ground-breaking by advancing new methodologies. Ironically, from the vantage point of (retrospective) meta-analysis, the fact that an innovative study inspired other studies to follow suit means that an innovative study will in hindsight obscure its own central role in the exploration of a phenomenon.  

To make this more formal, let us revisit the stylized example from Section \ref{Sec:Studies}\ref{Sec:ToyExample}. 
The introduction of new research methodology by the $11$\textsuperscript{th} study led to a drastic shift and widening of the collective posterior. In our hypothetical sequence of studies, this innovation also led future researchers to adopt the new methodology. This study was clearly highly influential. What do meta-analytic algorithms say about this study's importance?  

When performing meta-analysis, it is typical to look at the weights of studies (i.e., how much each study contributed to the pooled meta-analytic estimate) to determine their importance. Their weights are given by 
\begin{align}\nonumber
    \delta_i = \frac{w_i}{\sum_{i=1}^n w_i},
\end{align}
where $w_i = \frac{1}{\hat \sigma_i^2}$ for fixed-effect or $w_i = \frac{1}{\hat \sigma_i^2 +\hat \tau^2}$ for random effects meta-analysis. Here   $\hat \sigma_i$ are the standard errors reported by the studies and $\hat \tau$ is the meta-analysis estimate for effect heterogeneity \citep{Higgins_Thompson_2002}. 
To isolate the effect of innovation, our stylized example is set up in such a way that all studies' effects have the same standard error. We can thus immediately see that all studies will be accorded same weight -- even though study 11 was a pathbreaking and influential study, it appears as no different than any other study. Table \ref{tab:MAweights} shows the weights of the latest study when meta-analysis is performed sequentially. In stark contrast to our learning metric (see Figure  \ref{fig:syn_trace} (Right)), the meta-analysis weights imply diminishing returns.

Meta-analysis weights, even when meta-analysis is performed sequentially, fail to properly quantify the contribution of a methodologically innovative study. Innovation, as opposed to pure methodological diversity, is inherently sequential and the value of innovation cannot be captured purely by the corresponding studies' contribution to the aggregate effect estimate. By focusing on the effect a study has on the collective posterior at each point in time, which may drastically widen when innovation leads to a surprising outcome, our method of quantifying learning is able to pick up on the special role the $11$\textsuperscript{th} study plays in the stylized sequence of studies. 

While our discussion has centered on a stylized example, it has important implications for judging the value of studies in the real world. When meta-analysis weights are used to gauge a study's contribution to the literature, studies with small reported standard errors receive higher weights. This essentially favors large (and oftentimes expensive) studies while diminishing the importance of smaller and innovative pilot studies. Through the lens of our method for quantifying learning, the precision of a study is only one factor contributing to its impact. By being able to quantify the value of innovation, our method gives due credit to pioneers. 


\section{Lindley's Learning Metric}
\label{Sec:Lindley}
In his seminal paper, Lindley \citep{Lindley_1956} develops a measure of how much information a study has provided: 
\begin{align}
    \mathcal{I}(\pi(\vartheta),y) \, &= \, H(\pi(\vartheta\lvert y)) - H(\pi(\vartheta)) \\ \nonumber
    \, &:= \,\mathbb{E}_{\pi(\vartheta\lvert y)}\bigg [ \log\pi(\vartheta\lvert y) \bigg] - \mathbb{E}_{\pi(\vartheta)}\bigg [ \log\pi(\vartheta) \bigg]
\end{align}
This expression depends on the observed data $y$; some experimental outcomes are more informative than others. In addition to its information-theoretic appeal, Lindley's measure $I(\pi(\vartheta),y)$ can also be understood as the reduction of risk corresponding to a decision problem where one has to decide which distribution to report about an unknown quantity $\vartheta$ \citep{Parmigiani_1994}. The utility function $U(\vartheta, \varphi) = \log \varphi(\vartheta)$ implied by Lindley's measure is appealing from this decision-theoretic perspective as it induces decision makers to report their current beliefs in the form of their posterior. 


While theoretically appealing, Lindley's measure $I(\pi(\vartheta),y)$ can be hard to compute in practice \citep{Batu_Dasgupta_Kumar_Rubinfeld_2002}. 
Similar to the Wasserstein-2 distance, see Section \ref{Sec:LearningMetric}, Lindley's measure of information has a closed-form solution in the situation in which prior and posterior are normal. Let $\pi(\vartheta\lvert y) = \text{normal}(\mu_{post},\sigma_{post})$ and $\pi(\vartheta) = \text{normal}(\mu_{prior},\sigma_{prior})$, then 
\begin{align}
    \mathcal{I}(\pi(\vartheta),y) = \log{\sigma_{post}}-\log{\sigma_{prior}}.
\end{align}
This reveals why Lindley's information measure is unappealing in practice. Realistically, the sequential meta-analysis model does not coincide with the true data-generating process. In this situation, we are not guaranteed that the posterior mean concentrates on the true value of $\vartheta$ -- and different studies (for example by removing bias) may impact our shared knowledge by drastically changing the posterior mean. We thus recommend using the Wasserstein distances as learning metrics (Equation \ref{eq:LM}) as they are sensitive to changes in the entire posterior distribution.
